% -----------------------------------------------------------
\section{New Birth}
% -----------------------------------------------------------
	\label{NEWBIRTH}
	
	The phrase ``new birth'' doesn't occur in the scriptures, and neither do any of its possible permutations with other forms. It's just the idea of being born again and being changed. The scriptures call it so many different things that I'm just using the term ``new birth'' to keep track of them all.
	
% % ----------------------
% \subsection{See also}
% % ----------------------

% % ----------------------
% \subsection{1828 American Dictionary of the English Language} % *1828
% % ----------------------
	% \label{NEWBIRTH-1828}

	% \begin{compactitem}
		% \item
	% \end{compactitem}

% % ----------------------
% \subsection{Oxford English Dictionary} % *OED
% % ----------------------
	% \label{NEWBIRTH-OED}

	% \begin{compactitem}
		% \item[]
	% \end{compactitem}

% ----------------------
\subsection{Scriptural Usage} % *SCRIPTURES
% ----------------------
	\label{NEWBIRTH-SCRIPTURES}

	% % -----
	% \subsubsection{Old Testament} % *OT
	% % -----
		% \label{NEWBIRTH-OT}
	
		% % --
		% \paragraph{KJV}
		% % --
			% \label{NEWBIRTH-OT-KJV}
	
			% \begin{tabular}{ll}
				% Total occurrences in the OT & 
			% \end{tabular}
			
			% \begin{compactdesc}
				% \item[]
			% \end{compactdesc}
			
		% % --
		% \paragraph{Hebrew Bible}
		% % --
			% \label{NEWBIRTH-HEBREW}
		
			% %
			% \subparagraph{\texorpdfstring{\he{sample word}}{}}
			% %
				% \label{PAGA} % whatever is in step bible without periods
			
				% \begin{compactdesc}
					% \item[HALOT , PDF ] \textbf{}
						% \begin{compactitem}
							% \item[1]
						% \end{compactitem}
					% \item[BDB , PDF ] \textbf{}
						% \begin{compactitem}
							% \item[1]
						% \end{compactitem}
					% \item[DCH PDF ] \textbf{}
						% \begin{compactitem}
							% \item[1]
						% \end{compactitem}
				% \end{compactdesc}
				
				% \begin{compactdesc}
					% \item[Some OT uses] \textbf{}
						% \begin{compactdesc}
							% \item[]
						% \end{compactdesc}
				% \end{compactdesc}
			
		% % --
		% \paragraph{Septuagint}
		% % --
			% \label{NEWBIRTH-LXX}
		
			% %
			% \subparagraph{\texorpdfstring{\gr{sample word}}{}}
			% %
		
	% -----
	\subsubsection{New Testament} % *NT
	% -----
		\label{NEWBIRTH-NT}
	
		% --
		\paragraph{KJV}
		% --
			\label{NEWBIRTH-NT-KJV}
			
	
			% \begin{tabular}{ll}
				% Total occurrences in the NT & 
			% \end{tabular}
			
			\begin{compactdesc}
				\item[{\hyperref[MATTHEW-18-3]{Matthew 18:3--6}}]
				\item[{\hyperref[MATTHEW-19-13]{Matthew 19:13--15}}]
				\item[{\hyperref[MARK-8-1]{Mark 8:1--9}}] \phantom{.}
					\begin{compactitem}
						\item One way of viewing these stories about the food is to see them as a display of the abundance that Christ creates in us, where there was little before.
					\end{compactitem}
				\item[{\hyperref[LUKE-6-31]{Luke 6:31--35}}]
				\item[{\hyperref[JOHN-1-6]{John 1:6--13}}]
				\item[{\hyperref[JOHN-3-1]{John 3:1--21}}]
				\item[{\hyperref[JOHN-4-9]{John 4:9--15}}] \phantom{.}
					\begin{compactitem}
						\item See also \hyperref[JOHN-7-38]{John 7:38--39} here.
					\end{compactitem}
				\item[{\hyperref[JOHN-5-17]{John 5:17--47}}]
				\item[{\hyperref[JOHN-6]{John 6}}] \phantom{.}
					\begin{compactitem}
						\item This is Christ's discussion of the bread of life. A lot of what he says is relevant to the new birth.
					\end{compactitem}
				\item[{\hyperref[ROMANS-12-1]{Romans 12:1--2}}]
				\item[{\hyperref[ROMANS-6-3]{Romans 6:3--10}}] \phantom{.}
				\item[{\hyperref[2CORINTHIANS-3-6]{2 Corinthians 3:6}}]
				\item[{\hyperref[2CORINTHIANS-3-18]{2 Corinthians 3:18}}] \phantom{.}
					\begin{compactitem}
						\item Make sure to check out my note about ``beholding as in a glass,'' which creates some connections to Alma 5.
					\end{compactitem}
				\item[{\hyperref[2CORINTHIANS-4-10]{2 Corinthians 4:10--12}}]
				\item[{\hyperref[2CORINTHIANS-4-16]{2 Corinthians 4:16}}]
				\item[{\hyperref[2CORINTHIANS-5-17]{2 Corinthians 5:17--21}}]
				\item[{\hyperref[GALATIANS-4-1]{Galatians 4:1--7}}]
				\item[{\hyperref[GALATIANS-4-29]{Galatians 4:29}}]
				\item[{\hyperref[GALATIANS-6-15]{Galatians 6:15}}]
				\item[{\hyperref[EPHESIANS-2-1]{Ephesians 2:1--10}}]
				\item[{\hyperref[EPHESIANS-2-13]{Ephesians 2:13--22}}] \phantom{.}
					\begin{compactitem}
						\item Probably just all over chapter 2 really.
					\end{compactitem}
				\item[{\hyperref[EPHESIANS-3-14]{Ephesians 3:14--21}}] \phantom{.}
					\begin{compactitem}
						\item Especially verse 16, but other verses too.
					\end{compactitem} 
				\item[{\hyperref[EPHESIANS-4-17]{Ephesians 4:17--32}}] \phantom{.}
					\begin{compactitem}
						\item The verses at the end of this passage describe things we should be doing \textit{after} the new birth.
					\end{compactitem}
				\item[{\hyperref[PHILIPPIANS-3-20]{Philippians 3:20--21}}]
				\item[{\hyperref[COLOSSIANS-1-9]{Colossians 1:9--29}}]
			\end{compactdesc}
			
		% % --
		% \paragraph{Greek New Testament} % *GREEK
		% % --
			% \label{NEWBIRTH-GREEK}
		
			% %
			% \subparagraph{\texorpdfstring{\gr{sample word}}{}}
			% %
				% \label{KATALLAGE}
			
				% \begin{compactdesc}
					% \item[BDAG , PDF ] \textbf{}
						% \begin{compactitem}
							% \item 
						% \end{compactitem}
					% \item[LSJ , PDF ] \textbf{}
						% \begin{compactitem}
							% \item 
						% \end{compactitem}
				% \end{compactdesc}
				
				% \begin{compactdesc}
					% \item[Some NT uses] \textbf{}
						% \begin{compactdesc}
							% \item[]
						% \end{compactdesc}
				% \end{compactdesc}
		
	% -----
	\subsubsection{Book of Mormon} % *BOM
	% -----
		\label{NEWBIRTH-BOM}
	
		% \begin{tabular}{ll}
			% Total occurrences in the BOM & 
		% \end{tabular}
		
		\begin{compactdesc}
			\item[{\hyperref[MOSIAH-5-2]{Mosiah 5:2--15}}]
			\item[{\hyperref[MOSIAH-18-7]{Mosiah 18:7--22}}]
			\item[{\hyperref[MOSIAH-27-23]{Mosiah 27:23--31}}]
			\item[{\hyperref[MOSIAH-28-1]{Mosiah 28:1--4}}]
			\item[{\hyperref[ALMA-5]{Alma 5}}] \phantom{.}
				\begin{compactitem}
					\item Many verses in this chapter reference the new birth.
					\item They go as far into it as 54, which is an interesting verse too.
				\end{compactitem}
			\item[{\hyperref[ALMA-7-14]{Alma 7:14--16}}]
			\item[{\hyperref[ALMA-13]{Alma 13}}]
			\item[{\hyperref[ALMA-18-41]{Alma 18:41--19:36}}] \phantom{.}
				\begin{compactitem}
					\item Lamoni's experience.
				\end{compactitem}
			\item[{\hyperref[ALMA-22-15]{Alma 22:15--27}}]
			\item[{\hyperref[ALMA-26-3]{Alma 26:3}}]
			\item[{\hyperref[ALMA-26-14]{Alma 26:14--15}}]
			\item[{\hyperref[ALMA-36-6]{Alma 36:6--26}}]
			\item[{\hyperref[MORMON-9-6]{Mormon 9:6}}]
			\item[{\hyperref[MORONI-6-1]{Moroni 6:1--4}}]
			\item[{\hyperref[MORONI-8-10]{Moroni 8:10--}}] \phantom{.}
				\begin{compactitem}
					\item Note what it means to be ``alive in Christ.''
					\item See verse 19 as well, and 22
				\end{compactitem}
		\end{compactdesc}
		
	% -----
	\subsubsection{Doctrine and Covenants} % *DC
	% -----
		\label{NEWBIRTH-DC}
	
		% \begin{tabular}{ll}
			% Total occurrences in the DC & 
		% \end{tabular}
		
		\begin{compactdesc}
			\item[{\hyperref[DC60-7]{DC 60:7}}]
			\item[{\hyperref[DC61-33]{DC 61:33--34}}]
			\item[{\hyperref[DC84-33]{DC 84:33}}]
		\end{compactdesc}
		
	% % -----
	% \subsubsection{Pearl of Great Price} % *PGP
	% % -----
		% \label{NEWBIRTH-PGP}
	
		% \begin{tabular}{ll}
			% Total occurrences in the PGP & 
		% \end{tabular}
		
		% \begin{compactdesc}
			% \item[]
		% \end{compactdesc}
		
% % ----------------------
% \subsection{Key Phrases} % *KEYPHRASES *PHRASES
% % ----------------------
	% \label{NEWBIRTH-PHRASES}

	% \begin{compactdesc}
		% \item[]
	% \end{compactdesc}
	
% % ----------------------
% \subsection{Bible Dictionaries} % *DICTIONARIES
% % ----------------------
	% \label{NEWBIRTH-BD}

% % ----------------------
% \subsection{Restoration Usage} % *RESTORATION
% % ----------------------
	% \label{NEWBIRTH-RESTORATION}

	% % -----
	% \subsubsection{Joseph Smith} % *JS
	% % -----
		% \label{NEWBIRTH-JS}
	
		% \begin{compactdesc}
			% \item[DATE/SOURCE]
		% \end{compactdesc}
	
	% % -----
	% \subsubsection{Brigham Young} % *BY
	% % -----
		% \label{NEWBIRTH-BY}
	
		% \begin{compactdesc}
			% \item[DATE/SOURCE]
		% \end{compactdesc}
	
% ----------------------
\subsection{Commentary and Studies} % *COMMENTARY *STUDIES
% ----------------------
	\label{NEWBIRTH-COMMENTARY}

	% -----
	\subsubsection{Key Points}
	% -----
		\label{NEWBIRTH-KEYS}
	
		\begin{compactitem}
			\item Patterns characteristic of the new birth
				\begin{compactitem}
					\item Being able to show forth wonders (\hyperref[MATTHEW-14-2]{Matthew 14:2}).
					\item A clear vision of yourself as evil, wicked, or otherwise unworthy (\hyperref[MOSIAH-4-2]{Mosiah 4:2}).
					\item Crying out to Christ and/or God for mercy (\hyperref[MOSIAH-4-2]{Mosiah 4:2}).
					\item A declaration of faith in Christ and/or God (\hyperref[MOSIAH-4-2]{Mosiah 4:2}).
					\item A visitation of the Spirit (\hyperref[MOSIAH-4-3]{Mosiah 4:3}, \hyperref[MOSIAH-5-3]{5:3}).
					\item Receiving a remission of sin (\hyperref[MOSIAH-4-3]{Mosiah 4:3}); see \hyperref[REMISSION]{Remission} for more.
					\item Feelings of joy, peace, and other positive things (\hyperref[MOSIAH-4-3]{Mosiah 4:3}).
					\item A changed heart (\hyperref[MOSIAH-5-2]{Mosiah 5:2}, \hyperref[MOSIAH-5-7]{5:7}, \hyperref[ALMA-5]{Alma 5}, \hyperref[ALMA-19-33]{Alma 19:33})
					\item An absence of the desire to do evil or wrong (\hyperref[MOSIAH-5-2]{Mosiah 5:2}, \hyperref[ALMA-19-33]{Alma 19:33}, \hyperref[ALMA-22-15]{Alma 22:15})
				\end{compactitem}
		\end{compactitem}
		
	% -----
	\subsubsection{Becoming the sons/children of God}
	% -----
		\label{NEWBIRTH-sons}
		
		\begin{compactitem}
			\item (See the discussion of the ``\hyperref[NEWBIRTH-children-christ]{children of Christ}'' below.)
			\item \hyperref[DEUTERONOMY-32-5]{Deuteronomy 32:5}
			\item \hyperref[ROMANS-8-14]{Romans 8:14--17}
			\item \hyperref[2CORINTHIANS-6-18]{2 Corinthians 6:18}
			\item \hyperref[GALATIANS-3-26]{Galatians 3:26--29}
			\item \hyperref[GALATIANS-4-1]{Galatians 4:1--7}
			\item \hyperref[PHILIPPIANS-2-15]{Philippians 2:15}
			\item \hyperref[MOSIAH-5-2]{Mosiah 5:2--15}
			\item \hyperref[MOSIAH-18-22]{Mosiah 18:22}
			\item \hyperref[3NEPHI-9-17]{3 Nephi 9:17}
				\begin{compactitem}
					\item Also verse 20.
				\end{compactitem} 
			\item \hyperref[3NEPHI-12-44]{3 Nephi 12:44--45}
			\item \hyperref[MORONI-7-48]{Moroni 7:48}
			\item \hyperref[DC11-30]{DC 11:30}
			\item \hyperref[DC25-1]{DC 25:1} \phantom{.}
				\begin{compactitem}
					\item Note the older versions of this verse, where this bit about being ``sons and daughters'' in the kingdom doesn't appear.
				\end{compactitem}
			\item \hyperref[DC34-3]{DC 34:3} \phantom{.}
				\begin{compactitem}
					\item This is an interesting verse because it suggests that, in sometimes referring to people as ``sons'' or ``daughters,'' Christ is speaking of them in their capacity as people who have been changed by the Atonement or who have adopted this new relationship with him.
				\end{compactitem}
			\item \hyperref[DC35-2]{DC 35:2} \phantom{.}
				\begin{compactitem}
					\item Note that this also speaks to the priesthood---we acquire more power in it by being able to adopt this in-dwelling relationship with Christ (who in turn has an in-dwelling relationship with the Father).
					\item See my notes at 35:2; this verse is fundamental in understanding this extremely important connection.
				\end{compactitem}
			\item \hyperref[DC39-4]{DC 39:4--5}
			\item \hyperref[DC41-6]{DC 41:6}
			\item \hyperref[DC42-52]{DC 42:52}
			\item \hyperref[DC45-8]{DC 45:8}
			\item \hyperref[DC50-43]{DC 50:43}
		\end{compactitem}
		
	% -----
	\subsubsection{Becoming like a child}
	% -----
		\label{NEWBIRTH-child}
		
		% --
		\paragraph{Mosiah 3:19}
		% --
		
			\hyperref[MOSIAH-3-19]{Mosiah 3:19}, among other scriptures, instructs us to become ``as a child'' in our effort to reconcile ourselves with God. Following that instruction is a list of qualities which apparently expand on the idea of what it is to be as a child (though see \hyperref[MOSIAH-3-19-NOTE-child]{here} for a different reading). The qualities listed in that verse are:
			
				\begin{compactitem}
					\item submissive
					\item meek
					\item humble
					\item patient
					\item full of love
					\item willing to submit to what the parents ``inflict'' on him
				\end{compactitem}
				
			Only recently (late 2020) has the strangeness of this list really hit me. It is strange in a number of ways. First, when you think of children, these are not at all the qualities that come immediately to mind. If I were to describe my own children, the first things I would say are energetic, loud, happy, emotional, funny, strong-willed, silly, curious, growing, learning, and so on. Of the list in Mosiah 3:19 the only one which might make my list would be ``full of love.'' So this list is strange because the qualities are not the ones most people associate with children.
			
			Another reason it's strange is that it's simply false most of the time that my kids have the qualities on this list. My kids are \textit{not} submissive to what I want, at least most of the time, and they are \textit{not} patient really at all. Meekness and humility seem to be a little different. When I think about my kids, I don't think the words ``meek'' and ``humble'' apply to them at all, not in the sense that they don't have those qualities, but in that the terms just are not appropriate for describing children in general; it's almost like a category mistake. I'm not sure I want meek children, in any case. I'd much rather they have some fire in them which they learn to temper as they grown, than that they're just mild and gentle all the time. So this list is strange because it's false in general that at least some of these qualities are characteristic of children. They aren't.
			
			So what's going on with this part of the verse then? There is the \hyperref[MOSIAH-3-19-NOTE-child]{alternative reading} of the qualities, which may explain it. It could also be that there is a way of being as a child \textit{when you are actually an adult} that is different from being as a child when you are actually a child. In the former case a person has many years of experience that children lack, and so is much more knowledgeable about the way the world works, and will (possibly) have a long history of interactions with the Lord. In those cases (to me at least) it would make more sense to use terms like ``meek'' and ``humble'' and ``patient'' to describe that person, and in fact I can see why an adult would \textit{want} to develop those qualities, and how it would feel to them (an adult) that in developing those qualities they are becoming more childlike.
			
			Another point to mention is that there are times when children do show the qualities listed in the verse more, and that's when they are scared, hurt, or otherwise unsure of themselves. In situations like that I think they do then become submissive, are more willing to be patient, and maybe it even becomes more appropriate to describe them as being meek and humble. Seeing the verse like this is a useful way of understanding when the new birth can have its greatest impact on us, but we also stand to gain from being submissive to the will of the Lord even when we aren't scared or hurt.
			
			All this said, I still strongly believe (and this is something I have come to believe somewhat recently) that it is important to develop the other more energetic qualities as part of the new birth. They're going to make us better disciples of Christ, they're going to have an uplifting impact on those around us, and, perhaps most importantly of all, they are going to make \textit{us} happier. \textit{We ourselves} are going to enjoy our life more when we develop those qualities. \textit{We} are going to have more fun, look forward to things more, be more playful, more spontaneous, and all of this is going to lead to a life that is just more full of goodness and happiness.
			
			Finally, just as a side note, I do wonder about the ``inflict'' language at the end of the verse. The connotation of that word is fairly negative and so is the actual \hyperref[INFLICT-1828]{definition} in the 1828 dictionary, so I'm not totally sure what to think about that.
			
	% -----
	\subsubsection{Children of Christ}
	% -----
		\label{NEWBIRTH-children-christ}
		
		Mosiah 5:7 reads as follows:
		
			\begin{quote}
				And now because of the covenant which ye have made, ye shall be called the children of Christ, his sons and his daughters; for behold, this day he hath spiritually begotten you, for ye say that your hearts are changed through faith on his name; therefore ye are born of him and have become his sons and his daughters.
			\end{quote}
			
		Here the people are called the children of Christ ``\textit{because of the covenant}'' which they have made. The covenant is what gives them the title of ``children of Christ.''
		
		Compare just receiving the title of ``children of Christ'' to actually \textit{being} children of Christ---here is Moroni 7:19:
		
			\begin{quote}
				Wherefore I beseech of you, brethren, that ye should search diligently in the light of Christ that ye may know good from evil. And if ye will lay hold upon every good thing and condemn it not, ye certainly will be a child of Christ. 
			\end{quote}
			
		Here the text says we won't just be \textit{called} a child of Christ, we'll \textit{be} a child of Christ. The reason isn't the covenant but the action of ``laying hold upon every good thing.'' (2/22/2024: Actually I think this is incomplete---they become his sons and daughters at the end of Mosiah 5:7. It's not just being called that.)
		
		There is a similarity in temple ordinances. In the endowment people are anointed \textit{to become} kings and priests, queens and priestesses, but in the second anointing, people are anointed to \textit{be} those things.
		
		The phrase ``child of Christ'' also appears at \hyperref[4NEPHI-1:17]{4 Nephi 1:17} and \hyperref[MORMON-9-26]{Mormon 9:26}.
		
		See also:
			\begin{compactitem}
				\item The parts in \hyperref[MOSIAH-14]{Mosiah 14--15} about being the ``seed'' of Christ.
					\begin{compactitem}
						\item 15:11, for one
					\end{compactitem}
				\item \hyperref[ETHER-3-14]{Ether 3:14}
				\item \hyperref[DC93-22]{DC 93:22}
			\end{compactitem}